\section{Discussion and conclusions}
\label{sec:discussion}

The model connects a local routing approximation to finite-fleet intercity
scheduling. Wave-capacity tables feed a time-expanded MILP, whose vehicle
inventories and freight precedence constraints coordinate collection,
line-haul, distribution, direct service, and independent empty repositioning.
Rolling execution preserves the resulting whole-task commitments and their
inventory consequences across computational windows. A physical terminal
fleet target is a separate operating requirement.

The analytical results explain two different choices. Pooling vehicles on
the same local arc weakly increases modeled capacity, while the conditional
load formulas balance repeated stem travel against shipment lead time.
The collection and distribution loads differ because their delivery-time
exposure differs. In the stationary benchmark, a per-shipment cost minimizer
also determines an average-cost-optimal repeated wave among the stated
successive complete-wave policies. The full calendar-time argument and QVI
verification are retained in Appendix~\ref{app:dispatch}.

These interpretations depend on the declared approximations. Balanced
zoning and stop-equivalent demand require appropriate calibration for an
application; the aggregate direct duration must represent the vehicle's
complete task before it is credited to the destination inventory. The
stationary results additionally require common waves and separable batch
costs, whereas shared fleets, finite sorting, and independently timed
stages can couple batches. Table~\ref{tab:model-scope} keeps these analytical
assumptions distinct from the scheduling model. The implementation reports
partial execution when continuation fails and uses window-level solver
certificates with the scope stated in Appendix~\ref{app:replay}.
